% Irodov Part V §5.2 originals 76--86. Source PDF pp.208--210, printed 199--201.
\interferenceproblem{76}{199}
单色平面波垂直照射双缝，缝间距 $d=2.5$mm，屏距 $L=100$cm。若在一缝前插入厚 $h=10\,\mathrm{\mu m}$ 的平行玻璃片，干涉条纹偏移多少，向哪个方向？
\exsolution
玻璃片使该光路额外光程增加
$\Delta=(n-1)h$。
双缝的几何光程差为 $dx/L$，
中心亮纹移至满足
$dx/L=\Delta$ 的位置，
因此
\[
\boxed{|\delta x|=\frac{(n-1)hL}{d}}.
\]
取普通玻璃 $n\simeq1.50$：
\[
\boxed{|\delta x|=\frac{0.50(10^{-5})(1.0)}
{2.5\times10^{-3}}=2.0\,\mathrm{mm}}.
\]
图样\emph{向放玻璃片的狭缝一侧}偏移。
\excheck 单色光波长在平移量表达式中消掉，因为插片使所有条纹阶数共同平移；原题没有单独给玻璃折射率，数值按普通玻璃约定。

\interferenceproblem{77}{199--200}
图5.16双臂气体干涉仪中，两根相同玻璃管充空气、长度均
$\ell=10.0$cm，入射波长 $\lambda=589$nm。
把上管空气换成氨气后，条纹向上移动 $N=17$ 个间距。
已知空气折射率 $n=1.000277$，求氨气折射率。
\exsolution
一臂由 $n$ 变为 $n'$，额外光程差为
$\Delta=(n'-n)\ell$。
每经过 $\lambda$ 对应一个完整干涉条纹移位，
故
$N=(n'-n)\ell/\lambda$，
即
\[
\boxed{n'=n+\frac{N\lambda}{\ell}}
=1.000277+\frac{17(589\times10^{-9})}{0.10}
\simeq\boxed{1.000377}.
\]
\excheck 此为单程双臂干涉装置，若是米歇尔逊回程双程管才需额外因子2；原图5.16的气体管仅穿行一次。

\interferenceproblem{78}{200}
电磁波垂直入射到折射率分别为 $n_1,n_2$ 的两个各向同性非磁介质分界面。利用切向电、磁场连续条件证明：(a) 透射电场不发生相位跳变；(b) 从疏到密介质反射的电场相位跳变 $\pi$。
\exsolution
设入射、反射、透射电场有向复振幅为
$E_i,E_r,E_t$，
垂直入射时电场边界连续：
$E_i+E_r=E_t$。
磁场反射方向相反，而非磁介质中
$H=nE/Z_0$，故
$n_1(E_i-E_r)=n_2E_t$。
消去 $E_t$ 得菲涅耳振幅系数
\[
\boxed{\frac{E_r}{E_i}=\frac{n_1-n_2}{n_1+n_2}},
\qquad
\boxed{\frac{E_t}{E_i}=\frac{2n_1}{n_1+n_2}}.
\]
透射系数恒正，故透射电场与入射电场同相；
若 $n_2>n_1$，反射系数为负，反射电场相位增加 $\pi$。
\excheck 此结论针对电场的相位，反射磁场还须考虑传播方向反转造成的取向符号。

\interferenceproblem{79}{200}
白光平行光束以入射角 $\theta_1=52^\circ$ 照射
$n=1.33$ 的薄膜。要使反射光中波长 $\lambda=0.60\,\mathrm{\mu m}$ 的黄色最强，薄膜厚度应取何值？
\exsolution
假设膜的两侧为空气，只有一侧反射发生 $\pi$ 相位跳变。
在膜内折射角 $r$ 满足
$\sin r=\sin\theta_1/n$。
两反射光相长条件为
$2nd\cos r=(k+1/2)\lambda$，
所以
\[
\boxed{d=\frac{(2k+1)\lambda}
{4\sqrt{n^2-\sin^2\theta_1}}},\quad k=0,1,2,\ldots
\]
代 $n=1.33,\theta_1=52^\circ$，
$\boxed{d\simeq0.14(2k+1)\,\mathrm{\mu m}}$。
\excheck 若膜下方是比膜折射率更高的基底，两个反射相位翻转关系会改变；本题采用原书薄自由膜的通常边界条件。

\interferenceproblem{80}{200}
折射率为 $n=1.33$ 的薄膜，在入射角
$\theta=30^\circ$ 下，要求
$\lambda_1=0.64\,\mathrm{\mu m}$ 的光反射最强、
而 $\lambda_2=0.40\,\mathrm{\mu m}$ 的光反射恰极小。
求满足条件的最小膜厚。
\exsolution
记
$D=2d\sqrt{n^2-\sin^2\theta}$。
由一侧反射相位反转，
$\lambda_1$ 最大要求
$D=(m+1/2)\lambda_1$，
$\lambda_2$ 最小要求
$D=p\lambda_2$，
其中 $m\ge0,p\ge1$ 均为整数。
比较
$(m+1/2)(0.64)=0.40p$，
即 $8m+4=5p$，
最小非负整数解为
$m=2,p=4$。
故
\[
\boxed{d_{\min}=
\frac{(2+1/2)\lambda_1}
{2\sqrt{n^2-\sin^2\theta}}
=\frac{1.6\,\mathrm{\mu m}}
{2\sqrt{1.33^2-\sin^230^\circ}}
\simeq0.65\,\mathrm{\mu m}}.
\]
\excheck 本题是两个波长的\emph{同时}约束，不能分别对两个波长取各自的四分之一波厚度。

\interferenceproblem{81}{200}
为减小玻璃表面反射损耗，表面镀一层透明增透膜，
膜折射率取 $n'=\sqrt n$，其中 $n$ 为玻璃折射率。
垂直入射波长 $\lambda$ 的光时，两次反射振幅恰相等。
求使合反射强度为零的膜厚。
\exsolution
空气到膜、膜到玻璃的两次反射均发生 $\pi$ 相位翻转，
二者的\emph{相对反射附加相位}为零。
膜内往返光程差为 $2n'd$；
要让两反射波完全抵消，传播相位差须为奇数个 $\pi$：
\[
\frac{2\pi}{\lambda}2n'd=(2k+1)\pi.
\]
因此
\[
\boxed{d=\frac{(2k+1)\lambda}{4n'}
=\frac{(2k+1)\lambda}{4\sqrt n}},\quad k=0,1,\ldots
\]
\excheck 等反射幅度所需 $n'=\sqrt n$，不是一般写作 $n'=n$；源PDF数学抽取可能漏根号。

\interferenceproblem{82}{200}
波长 $\lambda=0.60\,\mathrm{\mu m}$ 的漫射单色光照在
$n=1.5$ 的薄膜上，反射光中入射角附近
$\theta=45^\circ$ 的相邻两个角度极大间隔为
$\Delta\theta=3.0^\circ$。
求膜厚。
\exsolution
等倾角反射条纹的光程差
$D(\theta)=2d\sqrt{n^2-\sin^2\theta}$。
相邻极大的阶数差一，故
$|\mathrm dD/\mathrm d\theta|\Delta\theta\simeq\lambda$。
求导：
$\mathrm dD/\mathrm d\theta
=-d\sin2\theta/\sqrt{n^2-\sin^2\theta}$，
因此
\[
\boxed{d\simeq
\frac{\lambda\sqrt{n^2-\sin^2\theta}}
{\sin2\theta\,\Delta\theta}
\simeq15\,\mathrm{\mu m}}.
\]
\excheck 角间隔必须先换成弧度：
$3^\circ=0.05236$rad；源书答复中根号可能因PDF文本抽取漏失。

\interferenceproblem{83}{200}
图5.17：单色光穿过屏孔后在距屏
$\ell$ 处的平行薄板反射，形成等倾角同心圆干涉环。
薄板厚度 $d$、折射率 $n$，
第 $i$ 与 $k$ 个暗环半径分别为
$r_i,r_k\ll\ell$。
求波长 $\lambda$。
\exsolution
若屏位于板前距离 $\ell$，反射光倾角 $\theta$ 的环落于屏上径向距离
$r\simeq2\ell\theta$（往返反射几何）。
膜内两束反射光的光程差为
$D(\theta)=2d\sqrt{n^2-\sin^2\theta}
\simeq2nd-(d/n)\theta^2$。
相隔 $|i-k|$ 个暗条纹阶数的光程差变化满足
$|D(\theta_i)-D(\theta_k)|=|i-k|\lambda$。
于是
\[
|i-k|\lambda\simeq\frac d n
\frac{|r_i^2-r_k^2|}{4\ell^2},
\quad
\boxed{\lambda\simeq
\frac{d\,|r_i^2-r_k^2|}
{4n\ell^2|i-k|}}.
\]
\excheck 这里的径向关系 $r\simeq2\ell\theta$ 来自反射光往返走过屏—板距离；若误写为 $\ell\theta$，结果会差四倍。

\interferenceproblem{84}{200}
波长 $\lambda$ 的平面光入射到两面夹角 $\alpha\ll1$ 的玻璃楔，入射面垂直楔棱，入射角为 $\theta_1$。求在垂直于\emph{反射光}的观察屏上相邻等厚明条纹间隔。
\exsolution
设楔膜厚沿表面坐标 $x$ 变化
$d(x)\simeq\alpha x$。
介质折射率为 $n$，
内部折射角满足
$n\sin r=\sin\theta_1$。
反射的光程差
$D(x)=2nd(x)\cos r
=2\alpha x\sqrt{n^2-\sin^2\theta_1}$，
相邻亮纹相差 $\lambda$，故楔面上的间隔
$\Delta x_s=\lambda/
[2\alpha\sqrt{n^2-\sin^2\theta_1}]$。
若屏面垂直于反射方向，沿屏面的图样距离多一个几何投影因子
$\cos\theta_1$：
\[
\boxed{\Delta x=
\frac{\lambda\cos\theta_1}
{2\alpha\sqrt{n^2-\sin^2\theta_1}}}.
\]
\excheck 区分楔面本身的条纹间隔与倾斜观察屏上的投影间隔。

\interferenceproblem{85}{201}
来自远处点源的波长
$\lambda=0.55\,\mathrm{\mu m}$ 的光垂直照射玻璃楔。
反射明条纹在楔面上间隔为
$\Delta x=0.21$mm；在离楔棱约
$\ell=1.5$cm 处条纹消失。
求：(a) 楔角；(b) 相对谱宽 $\Delta\lambda/\lambda$。
\exsolution
垂直入射时
$2nd(x)=2n\alpha x$。
相邻极大光程差变化 $\lambda$，
故 $\Delta x=\lambda/(2n\alpha)$；
取普通玻璃 $n\simeq1.50$：
\[
\boxed{\alpha=\frac{\lambda}{2n\Delta x}
\simeq8.7\times10^{-4}\,\mathrm{rad}
\simeq3'}.
\]
距楔棱 $\ell$ 处条纹秩序约
$m\simeq\ell/\Delta x$；
不同波长造成的阶数不确定性
$m\Delta\lambda/\lambda\sim1$，
故
\[
\boxed{\frac{\Delta\lambda}{\lambda}
\sim\frac{\Delta x}{\ell}
=\frac{0.21\,\mathrm{mm}}{15\,\mathrm{mm}}
\simeq0.014}.
\]
\excheck 相对谱宽是“条纹开始洗掉”的数量级估计，不是一个精确的线宽定义。

\interferenceproblem{86}{201}
平凸玻璃透镜的凸面与平玻璃板相接触，
凸面曲率半径 $R$、照明波长 $\lambda$，
求牛顿环在 $\Delta r\ll r$ 的远离中心区域，环宽
$\Delta r$ 随环半径 $r$ 的关系。
\exsolution
空气膜厚度在离中心 $r$ 处为
$h(r)\simeq r^2/(2R)$。
相邻\emph{明暗边界}对应膜往返光程差变化 $\lambda/2$，
故
$2[h(r+\Delta r)-h(r)]\simeq\lambda/2$。
近似
$h(r+\Delta r)-h(r)\simeq(r/R)\Delta r$，因此
\[
\boxed{\Delta r\simeq\frac{\lambda R}{4r}}
\quad\text{（相邻一明一暗条纹的间隔）}.
\]
若“相邻同类极大”定义为完整条纹周期，则距离为
$\boxed{\Delta r_{\rm bright\to bright}\simeq\lambda R/(2r)}$。
\excheck 原书书末采用前一种“环宽”定义，故为 $1/4$ 系数；两种计数方式不能混淆。
